\input{template.tex}

\ifPDFTeX
  \renewcommand{\rmdefault}{cmr}
\fi

\def\SheetNo{7}

\begin{document}

\sheettitle{Графы. Начало}

\begin{problems}

\item
\begin{parts}
  \item Докажите, что в графе число вершин нечётной степени чётно.
  \item Какое максимальное число рёбер может быть в графе с $n$ вершинами?
  \item Какое максимальное число рёбер может быть в графе с $n$ вершинами и $k$ компонентами связности?
\end{parts}

\item Постройте граф с $n$ вершинами, $m$ рёбрами и $k$ компонентами связности.
\par Здесь и далее «постройте граф с $n$ вершинами, $\ldots$» означает, что вы должны рассказать способ для любого $n$ построить искомый граф, либо рассказать, для каких $n$ такой граф существует и указать способ его построить, а для остальных $n$ доказать, что такого графа не существует. Аналогично следует поступить с другими параметрами, указанными в условии задачи.

\item Обозначим как $N(u)$ множество соседей вершины $u$, а как $N[u]$ множество, содержащее вершину $u$, а также соседей вершины $u$.
\begin{parts}
  \item Постройте граф с $n$ вершинами, в котором множества $N(u)$ совпадают для всех вершин $u$. Опишите все такие графы.
  \item Постройте граф с $n$ вершинами, в котором множества $N[u]$ совпадают для всех вершин $u$. Опишите все такие графы.
\end{parts}

\item Докажите, что для любого графа $G$ можно приписать каждой вершине $u$ число $d(u)$ так, что для любых вершин $u,v$ числа $d(u)$ и $d(v)$ имеют общий делитель, отличный от $1$, тогда и только тогда, когда в графе $G$ есть ребро $uv$.

\item В вершинах графа расставлены числа $\pm1$ таким образом, что для каждой вершины произведение чисел в вершинах, соседних с ней, отрицательно. Чему может быть равно произведение чисел во всех вершинах?

\item В графе $G$ $101$ вершина, а минимальная степень $\delta(G)\ge50$. Докажите, что $G$ связен.

\item Докажите, что если в графе $G$ с $n$ вершинами минимальная степень $\delta(G)>\frac{n-1}{2}$, то граф $G$ связен.

\item В связном графе нашлись два простых пути максимальной длины. Докажите, что у них есть хотя бы одна общая вершина.

\item Дан граф, содержащий $2n$ вершин и не менее чем $n^2+1$ ребро. Докажите, что в нём есть три вершины, попарно соединённые рёбрами.

\item В компании из $n$ человек среди любых четверых есть знакомый с остальными тремя. Докажите, что есть человек, который знает всех остальных.

\item В некоторой стране каждый город соединён с каждым дорогой с односторонним движением. Докажите, что найдётся город, из которого можно добраться в любой другой не более чем с одной пересадкой.

\item Пусть $T$ --- дерево с хотя бы $3$ вершинами. Докажите, что $T$ имеет хотя бы одну из двух следующих конфигураций: (*) два листа, смежных с одной и той же вершиной; (**) лист, смежный с вершиной степени $2$.

\item Вершины дерева $T$ покрашены в красный и синий цвет так, что любые две смежные вершины разноцветны, и красных не меньше, чем синих. Докажите, что у $T$ есть красный лист.

\item В выпуклом $n$-угольнике провели несколько диагоналей, не имеющих общих внутренних точек, которые разбили его на треугольники (это называется триангулировать многоугольник).
\begin{parts}
  \item Докажите, что получилось ровно $n-2$ треугольника.
  \item Пусть каждый треугольник --- вершина графа, а две вершины смежны, если треугольники имеют общую сторону. Докажите, что получилось дерево.
\end{parts}

\item Может ли сумма попарных расстояний между вершинами $25$-вершинного дерева быть равна $1225$?

\item В графе степень каждой вершины не меньше $1$. Докажите, что его можно разбить на $2$ части так, чтобы любая вершина была соединена с какой-то вершиной из другой части.

\item Множество клеток на клетчатой плоскости назовём ладейно связным, если из каждой его клетки можно попасть в любую другую, двигаясь по клеткам этого множества ходом ладьи (ладье разрешается перелетать через поля, не принадлежащие нашему множеству). Докажите, что ладейно связное множество из $100$ клеток можно разбить на пары клеток, лежащих в одной строке или в одном столбце.

\end{problems}

\end{document}
